% docs/introduccion.tex
\section{Introducción}

El diseño de mezclas de concreto estructural es una de las competencias fundamentales en la formación del ingeniero civil, y su dominio exige no solo el conocimiento de la normativa —como el método ACI\,211.1— sino también la capacidad de contrastar hipótesis con evidencia experimental propia, registrar resultados de manera sistemática y comparar ensayos de forma reproducible. Esta articulación entre teoría y práctica adquiere una dimensión crítica en el departamento del Meta (Colombia), donde las condiciones ambientales del territorio —temperaturas elevadas, humedad relativa alta y suelos con características particulares— inciden directamente sobre el desempeño del concreto y obligan a tomar decisiones de dosificación que no pueden reducirse a tablas genéricas.

En este contexto, los semilleros de investigación y los laboratorios de materiales de la Universidad Cooperativa de Colombia (UCC), sede Meta, representan el espacio natural donde los estudiantes construyen y validan ese conocimiento. Sin embargo, la actividad experimental del semillero requiere de una base digital que permita anclar sus hipótesis de diseño a datos reales: registrar con precisión las condiciones de cada ensayo, confrontar resultados entre sesiones y disponer de un repositorio organizado que soporte la discusión académica y el avance progresivo del proceso investigativo. Hoy esa base no existe como herramienta propia del grupo, lo que limita la trazabilidad de los experimentos, dificulta la comparación entre cohortes y hace que el conocimiento acumulado en cada práctica no quede debidamente sistematizado.

Frente a esa necesidad surge este trabajo de grado, que parte de una premisa clara: \textbf{el semillero necesita una herramienta cuyo núcleo sean los propios procesos experimentales}, de modo que los resultados obtenidos en laboratorio —y no estimaciones externas— fundamenten las recomendaciones de dosificación y sirvan de base para plantear y refutar hipótesis sobre el comportamiento del concreto bajo las condiciones ambientales del Meta. La herramienta que se propone, denominada \textbf{DIAPCE (Diseño y Análisis de Proporciones de Concreto Estructural)}, está dirigida exclusivamente a los estudiantes del semillero y del laboratorio de materiales de la UCC; no es un producto de uso industrial ni general. Su valor reside en convertir el acervo experimental del grupo en un sistema consultable, trazable y reproducible que potencie la calidad de la investigación formativa.

El documento se organiza de la siguiente manera: se expone primero el planteamiento del problema y la justificación que sustentan la necesidad de la herramienta; luego se presenta el marco teórico (normativa, fundamentos de dosificación y criterios de clasificación de aditivos); posteriormente se detalla la metodología y la arquitectura técnica de la solución; se describen las métricas de calidad y los resultados obtenidos; y finalmente se incluyen las conclusiones, recomendaciones y referencias bibliográficas.
